\documentclass[10pt,a4paper]{amsart}
\usepackage[margin=19mm]{geometry}
\usepackage{amsmath,amssymb,mathtools}
\usepackage[scr=rsfs,cal=euler]{mathalfa}
\usepackage{xfrac}
\usepackage[unicode,hidelinks,bookmarks=false]{hyperref}
\newcommand{\ie}{i.e.\ }
\newcommand{\cf}{cf.\ }
\newcommand{\resp}{respectively\ }
\newcommand{\Subsection}{\subsection}
\providecommand{\nobreakdash}{\nobreak\hbox{-}}
\setlength{\emergencystretch}{2em}
\multlinegap0pt
\usepackage{amssymb}
\usepackage{mathtools}
\usepackage{xcolor}
\newtheorem{proposition}{Proposition}
\newtheorem{definition}{Definition}
\newtheorem{lemma}{Lemma}
\newtheorem{theorem}{Theorem}
\newcommand{\C}{\mathbb{C}}
\newcommand{\GL}{\mathrm{GL}}
\DeclareMathOperator{\Mat}{Mat}      
\newcommand{\OO}{\mathcal{O}}
\newcommand{\ZZ}{\mathbb Z}
\newcommand{\fol}{\mathcal{F}} 
\title{A Poisson--Poincar\'e--Dulac for Poisson Connections}
\author{Maur\'icio Corr\^ea, Miguel Rodr\'iguez Pe\~na}
\date{Source: arXiv:2602.14132v2 (CC BY 4.0)}
\begin{document}
\maketitle

\noindent\emph{Source and licence.} A Poisson--Poincar\'e--Dulac for Poisson Connections. Version arXiv:2602.14132v2 (CC BY 4.0), distributed by arXiv under the Creative Commons Attribution 4.0 International licence, http://creativecommons.org/licenses/by/4.0/, as recorded in the arXiv licence metadata for this exact version and shown on the abstract page https://arxiv.org/abs/2602.14132v2. Full licence notice: \texttt{licenses/arxiv-2602.14132-CC-BY-4.0.txt}.\par\smallskip

\noindent\textit{Editorial scope.} Counted unit: the article's theorem thm:PPD-local (source0.tex lines 996--1003). The article's complete original proof of it is delivered with it (source0.tex lines 1005--1059, one \texttt{proof} environment of 55 lines). No mathematics is written, repaired or re-derived: the statement and the proof are the article's own lines, in the article's own notation, and no retained line is substituted or re-flowed.  Retained uncounted as the source-local context the proof needs: lines 618--630; lines 666--670; lines 695--705; lines 766--788; lines 921--933.  Nothing was omitted from any retained span, so the package carries no omission record.  No retained line contains a citation, so the wrapper carries no bibliography and no bibliography entry.  No retained line contains a figure environment, an image-inclusion command or drawing code, so no image is delivered and the package declares no asset dependency.  Editorial formatting only, in addition: an AMS \texttt{amsart} preamble that restates the article's own declarations and macros used by the retained lines, namely the environments \texttt{proposition}, \texttt{definition}, \texttt{lemma}, \texttt{theorem} on separate counters, together with the standard packages those lines need.  The retained environments print in this document as Proposition 1, Definition 1, Lemma 1, Theorem 1 in order of appearance, whereas the article prints the same items with the numbering of its own full document.  The complete native source of the article as distributed in the arXiv e-print for this exact version is bundled unchanged as \texttt{sources/194687/source0.tex}.
\begin{proposition}[Flatness of Euler--Poisson principal parts]\label{prop:EP-flat}
Let
$$
\Theta_0=\sum_{i=1}^r A_iX_i,
\qquad
X_i=\sigma^\# \Big(\frac{dz_i}{z_i}\Big),
$$
with constant matrices $A_i\in\Mat_{e\times e}(\C)$. If the residue tuple commutes, i.e.
$$
[A_i,A_j]=0\qquad\text{for all }i,j,
$$
then $\Theta_0$ is Poisson-flat. Conversely, if the bivectors $X_i\wedge X_j$, $1\le i<j\le r$, are linearly independent outside an analytic subset of codimension at least two, then flatness of $\Theta_0$ forces $[A_i,A_j]=0$ for all $i,j$.
\end{proposition}

\begin{equation}\label{eq:Theta-split}
\Theta=\Theta_0+V,
\qquad
\Theta_0:=\sum_{i=1}^r A_i^\circ\,\sigma^\#\Big(\frac{dz_i}{z_i}\Big)=\sum_{i=1}^r A_i^\circ X_i,
\end{equation}

\begin{definition}[Fixed Euler--Poisson principal part]\label{def:fixed-EP-principal-part}
Let $\Theta_0=\sum_i A_i^\circ X_i$ be as in \eqref{eq:Theta-split}.  We say that a logarithmic Poisson connection matrix $\Theta$ has \emph{fixed Euler--Poisson principal part} $\Theta_0$ if
\[
\Theta=\Theta_0+V
\]
and the remainder $V$ is residue-free and positive with respect to the boundary ideal.  Concretely, in the chosen Euler--Poisson chart this means that the component of $V$ in each logarithmic meridional generator $X_i$ has vanishing residue along $D_i$ and that, after the standard local residue-free normalization, one has
\begin{equation}\label{eq:residue-free-positive}
V\in I_D\cdot \Mat_{e\times e}\big(\fol_\sigma^{\log}(U)\big).
\end{equation}
Equivalently, $\Theta$ and $\Theta_0$ have the same Euler--Poisson residues along all local branches of $D$, and the remaining part begins in strictly positive boundary order.  If the generators $X_i$ satisfy relations, the condition is imposed modulo their effective boundary relations, as made precise by the Euler--Poisson quotient introduced below; the residue tuple is then understood only through the induced effective principal part.
\end{definition}

\begin{definition}\label{def:nonres}
The Euler--Poisson principal part $\Theta_0$ is effectively
Poisson--Poincar\'e--Dulac non-resonant at $p$ if, for every
$\kappa\neq\kappa'$ and every non-zero $\alpha\in\ZZ_{\ge0}^r$,
\[
        [\alpha]_\sigma
        +
        \bar\lambda_\kappa
        -
        \bar\lambda_{\kappa'}
        \neq
        0
        \quad\text{in }\mathfrak t_{\sigma,p}.
\]
For the analytic normalization theorem we also require the standard Poincar\'e
lower bound for these non-zero effective vectors, as needed in the majorant
argument.  In the independent-generator case $\mathfrak t_{\sigma,p}=\C^r$, the
condition is
\[
        \alpha+\lambda_\kappa-\lambda_{\kappa'}\neq0
        \quad\text{in }\C^r .
\]
\end{definition}

\begin{lemma}\label{lem:EP-homological-equation}
Let $R^{(N)}$ be the order $N$ defect term produced in the $I_D$-adic
normalization of a flat logarithmic Poisson connection with fixed
Euler--Poisson principal part $\Theta_0$.  Suppose that $\Theta_0$ satisfies
Definition~\ref{def:nonres}.  Then, after shrinking the adapted chart if
necessary, there exists $K^{(N)}\in\Mat_{e\times e}(I_D^N)$ such that
\[
        \delta K^{(N)}+[\Theta_0,K^{(N)}]
        \equiv
        -R^{(N)}
        \pmod {I_D^{N+1}} .
\]
\end{lemma}

\begin{theorem} \label{thm:PPD-local}
Let $p\in D\cap X^\circ$, and work in an adapted Euler--Poisson chart as above.  Assume that $\nabla^{\mathrm{Pois}}$ is Poisson-flat, that $\Theta$ has fixed Euler--Poisson principal part $\Theta_0=\sum_i A_i^\circ X_i$ in the sense of Definition~\ref{def:fixed-EP-principal-part}, and that the Euler--Poisson principal part is effectively PPD non-resonant at $p$ in the sense of Definition~\ref{def:nonres}.  Then, after shrinking $U$ around $p$, there exists a holomorphic gauge $H\in\GL_e(\OO_U)$ such that
\begin{equation}\label{eq:PPD-Theta}
\widetilde\Theta=(\delta H)H^{-1}+H\Theta H^{-1}
=\sum_{i=1}^r \widetilde A_i\,\sigma^\#\Big(\frac{dz_i}{z_i}\Big),
\end{equation}
where the matrices $\widetilde A_i\in\Mat_{e\times e}(\C)$ are constant and commuting.  The normal form is written relative to the chosen Euler--Poisson generators $X_i$; if these generators satisfy relations, the residue tuple is understood in the effective quotient $\mathfrak t_{\sigma,p}$.
\end{theorem}

\begin{proof}
We give the proof in the chosen effective PPD chart.  Since \(\Theta\) has fixed Euler--Poisson principal part, write
\[
\Theta=\Theta_0+V,
\qquad
\Theta_0=\sum_{i=1}^r A_i^\circ X_i,
\qquad
V\in I_D\cdot\Mat_{e\times e}(\fol_\sigma^{\log}(U)).
\]
The matrices \(A_i^\circ\) commute, hence \(\Theta_0\) is Poisson-flat by Proposition~\ref{prop:EP-flat}.  Put
\[
L(K):=\delta K+[\Theta_0,K]
\]
for the linearized homological operator.
We construct a holomorphic gauge
\[
H=I+\sum_{N\ge1}K^{(N)},
\qquad K^{(N)}\in\Mat_{e\times e}(I_D^N),
\]
by induction on the boundary order.  Suppose that \(H^{[N-1]}\) has already been constructed so that the
transformed connection is congruent to \(\Theta_0\) modulo \(I_D^N\), and let
\[
R^{(N)}:=(\delta H^{[N-1]})(H^{[N-1]})^{-1}
       +H^{[N-1]}\Theta(H^{[N-1]})^{-1}-\Theta_0
\]
be its order \(N\) defect.  Then
\(R^{(N)}\in\Mat_{e\times e}(I_D^N\fol_\sigma^{\log})\).  Flatness of both \(\Theta\) and \(\Theta_0\), together
with the induction hypothesis, gives the linearized Maurer--Cartan compatibility condition
\[
        \delta R^{(N)}+[\Theta_0,R^{(N)}]\equiv0
        \pmod {I_D^{N+1}}.
\]
Moreover the induction is performed inside the fixed-principal-part subspace, so the order \(N\) defect is
residue-free in the logarithmic meridional directions.
By Lemma~\ref{lem:EP-homological-equation}, the homological equation
\[
        L(K^{(N)})\equiv -R^{(N)}\pmod{I_D^{N+1}}
\]
has a solution with \(K^{(N)}\in\Mat_{e\times e}(I_D^N)\).  Taking
\[
H^{[N]}=(I+K^{(N)})H^{[N-1]}
\]
therefore improves the normalization by one boundary order.  The quadratic terms in \(K^{(N)}\) lie in
\(I_D^{2N}\subset I_D^{N+1}\) for \(N\ge1\), so no further order \(N\) contribution appears.

This produces a formal gauge.  The lower-bound clause in the effective non-resonance condition gives the
standard Poincar\'e majorant estimate for the off-diagonal homological inverses, while the diagonal primitives
are obtained by the usual holomorphic homotopy operators on polydiscs and preserve the \(I_D\)-order by
Cauchy estimates.  Hence the formal product converges on a smaller polydisc.  The limiting gauge
\(H\in\GL_e(\OO_U)\) transforms \(\Theta\) into a pure Euler--Poisson normal form
\[
        \widetilde\Theta=\sum_{i=1}^r\widetilde A_iX_i,
\]
with constant commuting residues; commutativity follows from flatness of the transformed principal part.
\end{proof}

\end{document}
